% Generated draft: requires mathematical and source review.
\documentclass[11pt]{article}
\usepackage[margin=25mm]{geometry}
\usepackage{fontspec}
\setmainfont{Latin Modern Roman}
\usepackage{amsmath,amssymb,mathtools,unicode-math}
\setmathfont{Latin Modern Math}
\usepackage{enumitem,needspace,xurl,hyperref}
\setlength{\parindent}{0pt}
\setlength{\parskip}{5pt}
\setlength{\emergencystretch}{3em}
\newcommand{\N}{\mathbb{N}}
\newcommand{\Z}{\mathbb{Z}}
\newcommand{\Q}{\mathbb{Q}}
\newcommand{\R}{\mathbb{R}}
\newcommand{\C}{\mathbb{C}}
\newcommand{\F}{\mathbb{F}}
\newcommand{\A}{\mathbb{A}}
\newcommand{\PP}{\mathbb{P}}
\newcommand{\E}{\mathbb{E}}
\newcommand{\eps}{\varepsilon}
\newcommand{\dd}{\,\mathrm{d}}
\begin{document}
\section*{UnsolvedMath ID: 30002041}
\textbf{Source problem code:} catalogue entry
\subsection{Definitions and mathematical statement}
Let \(X\) be an \(n\)-dimensional Gorenstein affine toric variety, and let \(\mathcal E=(E_0,\ldots,E_r)\) be a finite collection of rank-one reflexive \(\mathcal O_X\)-modules. Here reflexive means that the canonical map \(E_i\longrightarrow E_i^{\vee\vee}\) is an isomorphism, where \(E_i^\vee=\mathcal Hom_{\mathcal O_X}(E_i,\mathcal O_X)\); rank one means that the generic fiber has dimension one over the function field of \(X\). Set
\[
A=\operatorname{End}_{\mathcal O_X}\!\left(\bigoplus_{i=0}^{r}E_i\right).
\]
Assume that \(A\) is a consistent toric algebra in the sense intended by the supplied sources and that
\[
\operatorname{gldim}A=n=\dim X.
\]
The ground field and any restrictions on it are not specified in the supplied record.

The question is whether, for every such choice of \(X\) and \(\mathcal E\), the source's higher-quiver construction exists and its associated cellular complex, augmented to \(A\), is a minimal projective \((A,A)\)-bimodule resolution of \(A\), with the cellular form referred to as Theorem~0.2 in the original statement. A projective bimodule resolution means an exact augmented complex
\[
\cdots\longrightarrow P_2\longrightarrow P_1\longrightarrow P_0\longrightarrow A\longrightarrow 0
\]
in which every \(P_j\) is projective as an \((A,A)\)-bimodule. If \(k\) denotes the intended ground field, this is projectivity over \(A\otimes_k A^{\mathrm{op}}\), where \(A^{\mathrm{op}}\) is the opposite algebra.

\emph{Unresolved qualifications.} The supplied record does not reproduce the definitions of consistent toric algebra, higher quiver, its cellular complex, or the asserted form of Theorem~0.2. It also does not specify the convention for minimality in this setting. These terms retain their source-dependent meanings; the displayed general resolution does not specify the missing cellular terms or differentials. Consequently, the record supports the quantified question above, but does not supply enough information for a fully self-contained formulation of its cellular-resolution assertion.
\subsection{Short English statement}
For every endomorphism algebra of the stated rank-one reflexive sheaves on a Gorenstein affine toric variety, assuming consistency and global dimension equal to the dimension of the variety, must the higher quiver exist and encode a minimal projective bimodule resolution? The meanings of consistency, higher quiver, cellular form, and minimality require confirmation from the cited sources. The supplied registry labels this problem open; that status is not independently verified here.
\subsection{Sources}
\begin{itemize}
\item Local UnsolvedMath registry, supplied complete source record: permanent ID \(30002041\), title \emph{Cellular Resolutions for Consistent Toric Algebras}, source problem code \texttt{catalogue entry}, collection \emph{Oberwolfach Reports Open Problems}. The record gives proposed year \(2012\) and supplies no named proposer.
\item Supplied report citation: \emph{Toric Geometry (2013)}, Oberwolfach Reports, pp.~1231--1306, DOI \texttt{10.4171/owr/2012/21}. \url{https://doi.org/10.4171/owr/2012/21}. The supplied background describes Oberwolfach Reports as documenting mathematical workshops and their open-problem discussions. The report-year discrepancy is unresolved.
\item Craw--Quintero-V\'elez, \emph{Cellular resolutions of noncommutative toric algebras}. Supplied literature-source link: \url{https://eprints.gla.ac.uk/41620/1/41620.pdf}. This authorship is attribution for the cited work, not an identification of the problem's proposer. The original problem refers to Theorem~0.2, whose statement and precise location are not reproduced in the record.
\item Supplied literature assessment, checked on \texttt{2026-08-21}: the curation labels the problem open and reports that Craw--Quintero-V\'elez prove the cellular resolution for specified families in dimensions three and four, leave the extension to all consistent toric algebras as a conjecture, and that no general proof was found in that review. Its listed sources are the paper and report linked above. This is supplied dated curation derived from the OWR record, not an independently verified literature assessment; its status and cited claims require independent verification.
\item The registry labels the statement \texttt{corrected\_verified} and reports that the official report supplies previously omitted geometric and algebraic data and states the higher-quiver and cellular-resolution questions. This is supplied verification metadata, not a verification performed for this draft.
\end{itemize}
\end{document}
